% ---
% id: receiver/rem:res5why
% kind: rem
% title: "What the fifth residual checks that the others do not"
% summary: >-
%   The fifth residual is the only one that looks at the emitted photon-number
%   blocks, which can fail Hermiticity or positivity at the last step while the
%   other residuals pass. It is neither a repair nor an enclosure.
% prominence: standard
% status: accepted
% uses: [def:residuals]
% mentions: [thm:exact-and-evaluated, req:residual-refusal]
% assumes: []
% refs: []
% provenance: newly-derived
% lean:
% ---

\begin{remark}\label{rem:res5why}
Four of the other five residuals of Definition~\ref{def:residuals} check the
\emph{optics}: that the compiled matrix is a contraction, that the Grams are
positive, that the loss identity closes, that the POVM sums to the identity.
None of them looks at the object the artifact actually emits, the
photon-number block, and a block can fail Hermiticity or positivity at the last
floating-point step even when all four pass, because $r_4$ constrains only the
\emph{sum} over $S$ and $r_2$ lives on the Grams, one compilation stage
earlier. $r_5$ makes that defect visible where the receiver is in hand, with
the port, the setting and the photon number that produced it.

Two things $r_5$ is not. It is not a repair: a block that fails is refused and
never repaired (Requirement~\ref{req:residual-refusal}), because a repaired
block is a different operator and no claim of this contract holds for it. And
it is not an enclosure: $r_5=0$ says that the \emph{computed} block is
Hermitian and positive, not that it is close to the exact one, which is
Theorem~\ref{thm:exact-and-evaluated}\,(ii) exactly as for $r_1,\dots,r_4$.
\end{remark}
